%% ma: L11-B21-0014
%% lop: 11
%% chuong: 6
%% bai: 21
%% dang: 11.21.9
%% loai: TN4
%% muc: TH
%% dap_an: "B"
%% nguon: "(NBV)-ĐỀ SỐ 14-MỖI NGÀY 1 ĐỀ THI 2025.pdf (D14, MỖI NGÀY 1 ĐỀ - Đề số 14), Phần I Câu 10"
%% kien_thuc: "Bài 21 (phương trình mũ, lôgarit; ứng dụng thực tế)"
%% trang_thai: cho-duyet
%% ly_do: "Dạng mới 11.21.9 đề xuất, chờ giáo viên duyệt"
%% ghi_chu: "Đáp án gốc ghi C (9) sai: 8 ln(25/9) = 8,17 làm tròn là 8, đáp án đúng là B. Dùng hằng số e và lôgarit tự nhiên (Bài 19-21 lớp 11)."
\begin{ex}
Giả sử nhiệt độ của một loại đồ uống được xác định theo công thức $T=22+50e^{-\frac18t}$, $t\ge0$, trong đó $t$ (phút) là khoảng thời gian tính từ lúc pha chế đồ uống đó xong. Hỏi sau bao lâu kể từ lúc pha chế xong thì nhiệt độ của đồ uống đó là $40^\circ\mathrm C$? (Kết quả làm tròn đến hàng đơn vị.)
\choice{$7$}{\True $8$}{$9$}{$10$}
\loigiai{$22+50e^{-\frac t8}=40\Leftrightarrow e^{-\frac t8}=\dfrac{9}{25}\Leftrightarrow t=8\ln\dfrac{25}{9}\approx8{,}17$.
Làm tròn đến hàng đơn vị được $t\approx8$ phút.}
\end{ex}
